\section{Additional proof details}

\subsection{Weighted rank-one contraction}
\label{app:gram}

In unnormalized coordinates the weighted reflection is
$X^{\sharpop}=W_\eta^{-1}X^{\mathsf T}W_\eta$.  For
$q_p=u_pv_p^{\mathsf T}$,
\[
 q_q^{\sharpop}
 =W_\eta^{-1}v_qu_q^{\mathsf T}W_\eta.
\]
Consequently
\[
\begin{aligned}
\Tr(q_pq_q^{\sharpop})
&=\Tr\left(
u_pv_p^{\mathsf T}W_\eta^{-1}v_qu_q^{\mathsf T}W_\eta
\right)\\
&=(u_q^{\mathsf T}W_\eta u_p)
  (v_p^{\mathsf T}W_\eta^{-1}v_q).
\end{aligned}
\]
This also makes symmetry and nonnegativity of the divisibility
coefficients transparent.

For distinct primes, write $b=pqk$.  If $k$ contains $p$ or $q$, one
of $\mu(b/p),\mu(b/q)$ vanishes.  If any other prime divides $k$
twice, both vanish.  Otherwise
\[
 \mu(qk)\mu(pk)=\mu(k)^2=1.
\]
Thus
\[
\begin{aligned}
v_p^{\mathsf T}W_\eta^{-1}v_q
&=(pq)^{-2\eta}
\sum_{\substack{k\ {\rm squarefree}\\(k,pq)=1}}k^{-2\eta}\\
&=(pq)^{-2\eta}
\frac{C_\eta}
{(1+p^{-2\eta})(1+q^{-2\eta})}.
\end{aligned}
\]

\subsection{Uniform fourth-moment passage}
\label{app:fourier}

Let $F_N$ be increasing finite prime sets and
\[
 T_{s,N}=Z\left(\sum_{p\in F_N}p^{-s}E_p\right)Z^{-1}.
\]
On the critical line,
\[
 \sup_{t\in\mathbb R}
 \|T_{1/2+it}-T_{1/2+it,N}\|_4
 \le
 \|Z\|\,\|Z^{-1}\|
 \left(\sum_{p\notin F_N}p^{-2}\right)^{1/4}
 \longrightarrow0.
\]
The same estimate holds for the reflected leg, hence for $\cB$.
For $A,B\in\cS_4$, the telescoping identity and H\"older's trace
inequality give, on norm-bounded sets,
\[
 |\Tr(A^4)-\Tr(B^4)|
 \le C\|A-B\|_4.
\]
Therefore the finite trigonometric polynomials
$t\mapsto\Tr\cB_{1/2+it,N}^4$ converge uniformly.  Their Bohr Fourier
coefficients pass to the limit.

To isolate frequency $2\log(q/p)$, a term with indices
$(a,b,c,d)$ would require
\[
 \frac{bd}{ac}=\frac{q^2}{p^2}.
\]
Both sides are quotients of products of two primes.  Equality of
prime exponent vectors forces $a=c=p$ and $b=d=q$.  Thus the
coefficient computed in \cref{thm:fourth} is unique.

\subsection{Active commutant}

Suppose $KE_p=E_pK$ for an active coordinate $p$.  For $i\ne p$,
the $(i,p)$ entry of the left side is $K_{ip}$ and of the right side
is zero, so $K_{ip}=0$.  The $(p,i)$ entry similarly gives
$K_{pi}=0$.  Applying this to each active $p$ proves the active
one-dimensional blocks in \cref{thm:metric}.

If $K$ is positive and boundedly invertible, functional calculus
defines $K^{-1/2}$.  Then
\[
 U^*U=K^{-1/2}Z^*GZK^{-1/2}=I.
\]
Since $K$ is block diagonal with scalar active blocks, it commutes with
$D_s^A$.  This completes the unitary-collapse calculation without
assuming that the dormant block is diagonal.

\subsection{Modified atom product}

For one orthogonalized atom, the chiral block has eigenvalues
$\lambda_\pm=\pm p^{-1/2}$.  The third modified factors for
$I-zB$ are
\[
 (1-z\lambda_\pm)
 \exp\left(z\lambda_\pm+\frac{z^2\lambda_\pm^2}{2}\right).
\]
Multiplying the pair cancels the linear exponentials and yields
\[
 \left(1-\frac{z^2}{p}\right)e^{z^2/p}.
\]
Since
\[
 \log(1-z^2/p)+z^2/p=O(p^{-2})
\]
locally uniformly in $z$, the prime product converges normally.
